\begin{abstract}
Renderizar o contexto de um repositório como imagem reduz os tokens de entrada, mas não se sabe quanto isso custa em código que precisa passar em testes. Medimos esse custo em cinco modelos abertos de visão e linguagem, de 7B a 27B, em \nItemsAll itens do CoderUJB (Java) e do DevEval (Python), com o contexto em texto e em imagem a 10, 12 e 14 px. Os cinco modelos abertos passam menos testes com a imagem. O gap médio dos modelos, imagem menos texto, é \gapMeanUjb~pp \ci{\gapMeanUjbLo}{\gapMeanUjbHi} no CoderUJB e \gapMeanDev~pp \ci{\gapMeanDevLo}{\gapMeanDevHi} no DevEval. Por modelo, o gap vai de \gapMinUjb{} a \gapMaxUjb~pp e de \gapMinDev{} a \gapMaxDev~pp, e a razão de tokens de entrada vai de \tokenRatioMinAll{} a \tokenRatioMaxAll. Como resultado secundário, a penalidade varia entre itens (DP item:condição de \sdItemConditionUjb{} e \sdItemConditionDev{} em logit), mas no CoderUJB a comparação fora da amostra não apoia essa variação. Em uma extensão exploratória, o Claude Sonnet 5.5 e o Opus 5.5 também perdem, de \claudeLossMinAll{} a \claudeLossMaxAll~pp.
\end{abstract}
